\documentclass[10pt,a4paper]{amsart}
\usepackage[margin=19mm]{geometry}
\usepackage{amsmath,amssymb,mathtools}
\usepackage[scr=rsfs,cal=euler]{mathalfa}
\usepackage{xfrac}
\usepackage[unicode,hidelinks,bookmarks=false]{hyperref}
\newcommand{\ie}{i.e.\ }
\newcommand{\cf}{cf.\ }
\newcommand{\resp}{respectively\ }
\newcommand{\Subsection}{\subsection}
\providecommand{\nobreakdash}{\nobreak\hbox{-}}
\setlength{\emergencystretch}{2em}
\multlinegap0pt
\usepackage{xcolor}
\usepackage{amssymb}
\usepackage{mathtools}
\usepackage{enumerate}
\usepackage[normalem]{ulem}
\usepackage{float}
\usepackage{tikz}
\usepackage{dsfont}
\newtheorem{definition}{Definition}
\newtheorem{proposition}{Proposition}
\title{Quantum ergodicity in the Benjamini--Schramm limit for locally symmetric spaces}
\author{Farrell Brumley, Simon Marshall, Jasmin Matz, Carsten Peterson}
\date{Source: arXiv:2604.01075v1 (CC BY 4.0)}
\begin{document}
\maketitle

\noindent\emph{Source and licence.} Quantum ergodicity in the Benjamini--Schramm limit for locally symmetric spaces. Version arXiv:2604.01075v1 (CC BY 4.0), distributed by arXiv under the Creative Commons Attribution 4.0 International licence, http://creativecommons.org/licenses/by/4.0/, as recorded in the arXiv licence metadata for this exact version and shown on the abstract page https://arxiv.org/abs/2604.01075v1. Full licence notice: \texttt{licenses/arxiv-2604.01075-CC-BY-4.0.txt}.\par\smallskip

\noindent\textit{Editorial scope.} Counted unit: the article's proposition Sibd (source0.tex lines 1723--1726). The article's complete original proof of it is delivered with it (source0.tex lines 1728--1797, one \texttt{proof} environment of 70 lines). No mathematics is written, repaired or re-derived: the statement and the proof are the article's own lines, in the article's own notation, and no retained line is substituted or re-flowed.  Retained uncounted as the source-local context the proof needs: lines 1684--1689.  Nothing was omitted from any retained span, so the package carries no omission record.  No retained line contains a citation, so the wrapper carries no bibliography and no bibliography entry.  No retained line contains a figure environment, an image-inclusion command or drawing code, so no image is delivered and the package declares no asset dependency.  Editorial formatting only, in addition: an AMS \texttt{amsart} preamble that restates the article's own declarations and macros used by the retained lines, namely the environments \texttt{definition}, \texttt{proposition} on separate counters, together with the standard packages those lines need.  The retained environments print in this document as Definition 1, Proposition 1 in order of appearance, whereas the article prints the same items with the numbering of its own full document.  The complete native source of the article as distributed in the arXiv e-print for this exact version is bundled unchanged as \texttt{sources/197636/source0.tex}.
\begin{definition}\label{def:big}
Let $\Phi$ be a reduced root system. We say that a semistandard root subsystem $\Phi_0\subset\Phi$ is {\normalfont semi-dense} if, for every semistandard root subsystem $\Psi\subset\Phi$ we have
    \begin{equation}\label{eqn_root_ineq}
        |\Psi \cap \Phi_0| + \textnormal{rank}(\Psi) \geq \frac{1}{2} |\Psi|. 
    \end{equation}
    \end{definition}

\begin{proposition}
\label{Sibd}
    Suppose $\Phi$ is an irreducible root system of classical type. Then any extremal root subsystem is semi-dense. 
\end{proposition}

\begin{proof}
Let $\Phi_0\subset\Phi$ be extremal and let $\Psi\subset\Phi$ be semistandard. In fact, it will be enough to verify Definition \ref{def:big} with $\Psi$ standard. Indeed, writing $\Psi=w\Psi_{\rm st}$, where $\Psi_{\rm st}\subset\Phi$ is a standard root subsystem, we have $|\Psi \cap \Phi_0| = |\Psi_{\rm st} \cap w^{-1} \Phi_0|$. Since $\textnormal{rank}(\Psi)=\textnormal{rank}(\Psi_{\rm st})$ and $|w^{-1}\Phi_0|=|\Phi_0|$, inequality \eqref{eqn_root_ineq} is true for $\Psi$ if and only if it is true for $\Psi_{\rm st}$. 

We now let $\Psi$ be an arbitrary standard root subsystem of $\Phi$. To prove the proposition, it is enough to establish the inequality \eqref{eqn_root_ineq} for each of the irreducible factors of $\Psi$. We may therefore assume that $\Psi$ is irreducible. Moreover, we may assume that $\Psi\not\subset\Phi_0$ since otherwise the inequality \eqref{eqn_root_ineq}  is obviously true: $|\Psi|+\textnormal{rank}(\Psi)\geq |\Psi|/2$. We shall assume throughout the remainder of this proof that $\Psi$ is \textit{standard, irreducible and not contained in $\Phi_0$.}

Our strategy is to reduce  the verification of \eqref{eqn_root_ineq} to the base case $\Psi=\Phi$, which can be checked numerically. In this case we must verify
    \begin{equation}\label{basic-big-ineq}
        |\Phi_0|+\textnormal{rank}(\Phi) \geq \frac{1}{2} |\Phi|,
    \end{equation}
which can be checked by inspection of the following table:

    \begin{table}[H]
    \centering
 \begin{tabular}{|c|c|c|c|c|}
 \hline
  type of $\Phi$ & type of $\Phi_0$ & $|\Phi|$ & $|\Phi_0|$ \\ \hline  
  $A_n$ &$A_{n-1}$ & $n(n+1)$ & $(n-1)n$ \\
  $B_n$& $B_{n-1}$& $2 n^2$ & $2 (n-1)^2$  \\
  $C_n$& $C_{n-1}$& $2 n^2$& $2 (n-1)^2$ \\
  $D_n$& $D_{n-1}$& $2 n(n-1)$& $2(n-1)(n-2)$ \\
  \hline
 \end{tabular}
\end{table}
\noindent It remains then to show that when $\Psi\neq \Phi$ then $\Psi \cap \Phi_0$ is extremal in $\Psi$, in which case we can apply the base case to conclude. 
%(We will in fact encounter the degenerate setting when $\Psi$ is a singleton set, in which case $\Psi$ is of type $A_1$, $\Psi \cap \Phi_0=\emptyset$, and the inequality \eqref{eqn_root_ineq} is trivially verified, since $0 + 1 \geq \frac{1}{2} \cdot 2$.) 
It will be enough to prove this claim for $\Psi$ maximal. Note that maximal standard subsystems correspond to removing one simple root from $\Delta$.

We divide the proof of the above claim according to the root type.

\medskip

\noindent \textit{Type $A_n$.} The root system of type $A_n$ has a model as vectors in $\mathbb{R}^{n+1}$ of the form $e_i - e_j$ with $i \neq j$. Let $V$ be the subspace spanned by the roots (i.e., the hyperplane orthogonal to $e_1 + \dots + e_{n+1}$). A basis of simple roots is $e_1 - e_2, \dots, e_{n} - e_{n+1}$. Let $\Phi'_0$ be the extremal root subsystem generated by $e_2 - e_3, \dots, e_n - e_{n+1}$. This is exactly the intersection of $\Phi$ with the hyperplane orthogonal to $e_1$. All extremal root subsystems are of the form $\Phi_0=w \Phi'_0$ for some $w\in W$. The extremal root subsystems are therefore precisely the root subsystems of the form
\begin{equation}\label{eq:ext:root:sub:A}
\Phi_0 =\Phi\cap \{\langle e_j, x \rangle = 0\}
\end{equation} 
for some $1\leq j\leq n+1$.
%There are precisely $n+1$ of these each of the form $\Phi\cap \{\langle e_j, x \rangle = 0\}$ for some $1 \leq j \leq n+1$

Fix $1\leq j\leq n+1$ and suppose $\Phi_0$ is as in \eqref{eq:ext:root:sub:A}.
%    Suppose $\Phi_0= \Phi \cap \{ \langle e_j, x \rangle = 0 \}$. 
    We wish to show that $\Phi_0\cap \Psi$ is extremal in $\Psi$. Since $\Psi$ is a maximal standard root subsystem, there exists $1\leq l<n+1$ such that $\Psi$ is obtained by removing $e_l-e_{l+1}$ from the set of simple roots. 
    Let $V^{n+1}_l$ and $W^{n+1}_l$ denote the linear subspaces of $\mathbb{R}^{n+1}$ generated by $e_1,\ldots, e_l$ and $e_{l+1},\ldots,e_{n+1}$, respectively. Then $\Psi=(\Psi\cap V^{n+1}_l) \sqcup (\Psi\cap W^{n+1}_l)=:\Psi_1\sqcup\Psi_2$ with  $\Psi_1$ and $\Psi_2$ irreducible root systems of type $A_{l-1}$ and $A_{n-l}$, respectively. Moreover, $\Phi_0\cap \Psi$ equals either $\Psi_{1,0}\sqcup\Psi_2$ or $\Psi_1\sqcup \Psi_{2,0}$, depending on whether $j\le l$ or $j\ge l+1$, where $\Psi_{k,0}=\Psi_k\cap \{\langle e_j, x \rangle = 0\}$. By the characterization of extremal root subsystems given just before \eqref{eq:ext:root:sub:A} $\Psi_{k,0}$ is therefore extremal in $\Psi_k$, and thus $\Phi_0\cap \Psi$ is extremal in $\Psi$.
    
%\textcolor{gray}{    There are at most two simple roots in $\Phi$ which are not orthogonal to $e_j$, namely $e_{j-1} - e_j$ and $e_{j} - e_{j+1}$ (if $j = 1$ or $n+1$, then there is only one simple root not orthogonal to $e_j$). Since $\Psi$ is not contained in $\Phi_0$, it contains a simple root not orthogonal to $e_j$. Without loss of generality, we may suppose that $\Phi_0\cap \Psi$ is generated by the simple roots $e_k - e_{k+1}, \dots, e_{k+\ell - 1} - e_{k+\ell}$ with $k \leq j \leq k+\ell$ (and in particular is of type $A_\ell$). If $j = k$ or $k + \ell$, then it is clear that $\Phi_0\cap \Psi$ is extremal in $\Psi$. Otherwise, we can instead take as our basis for $\Psi$ the set $e_j - e_{k}, e_{k} - e_{k+1}, \dots, e_{j - 1} - e_{j+1}, e_{j+1} - e_{j+2}, \dots, e_{k + \ell - 1} - e_{k + \ell}$. Then $\Phi_0\cap \Psi$ is the root subsystem obtained by removing $e_j - e_k$ from the above basis. Thus the intersection is extremal (and in particular it is of type $A_{\ell-1}$).}

\medskip

    \noindent \textit{Types $B_n$ and $C_n$.} The $B_n$ and $C_n$ cases are virtually identical, and thus we just discuss $B_n$. The analysis is in turn very similar to the $A_n$ case. We have as a model of our root system vectors of the form $\pm e_i \pm e_j$ with $i \neq j$ and $\pm e_j$ inside of $\mathbb{R}^n$. We may take as basis $e_1 - e_2, e_2 - e_3, \dots, e_{n-1} - e_n, e_n$. Let $\Phi_0'$ be the extremal root subsystem generated by the $e_2 - e_3, \dots, e_{n-1} - e_n, e_n$. This is exactly the intersection of $\Phi$ with the hyperplane orthogonal to $e_1$. All other extremal root subsystems are of the form $\Phi_0=w \Phi_0'$. There are exactly $n$ of these all of the form $\Phi \cap \{\langle e_j, x \rangle = 0\}$ for some $1 \leq j \leq n$ (because the orbit of $e_1$ under $W$ is exactly vectors of the form $\pm e_j$).

    Suppose $\Phi_0 = \Phi \cap \{ \langle e_j, x \rangle = 0\}$. If $\Psi$ is obtained by removing one of the first $n-1$ simple roots, we again have $\Psi=(\Psi\cap V^{n}_l) \times (\Psi\sqcup W^{n}_l)=:\Psi_1\sqcup\Psi_2$ for some $1\le l<n$, but with $\Psi_2$ now of type $B_{n-l}$ or $C_{n-l}$. The claim then follows as in the $A_n$ case. If $\Psi$ is obtained by removing $e_n$, then $\Psi$ is the usual $A_{n-1}$ root system in $\mathbb{R}^{n}$ and we are also back in the previous case.
    
   %\textcolor{gray}{
    %If $j = 1$ then $e_1 - e_2$ is the only simple root not orthogonal to $e_j$. If $1 \leq j \leq n-1$, then $e_{j-1} - e_{j}$ and $e_{j} - e_{j+1}$ are the only simple roots not orthogonal to $e_j$. If $j = n$, then $e_{n-1} - e_n$ and $e_n$ are the only simple roots not orthogonal to $e_j$. Since $\Psi$ is not contained in $\Phi_0$, it contains a simple root not orthogonal to $e_j$. Thus, either $\Psi$ is of type $A_\ell$ in which case the analysis is exactly the same as before, or it is of type $B_\ell$, namely we have $e_k - e_{k+1}, \dots, e_{n-1} - e_n, e_n$. If $j \neq n$, then we can perform the same trick as before, i.e., choose the basis $e_j - e_k, e_k - e_{k+1}, \dots, e_{j-1} - e_{j+1}, \dots, e_{n-1} - e_n, e_n$ to conclude that the intersection is extremal (in particular of type $B_{\ell-1}$). Otherwise suppose $j = n$. Then we take as our basis: $e_n - e_k, e_k - e_{k+1}, \dots, e_{n-2} - e_{n-1}, e_{n-1}$ in which case we again see that the intersection is extremal.}

\medskip

\noindent \textit{Type $D_n$.} The type $D_n$ root system has a model as vectors in $\mathbb{R}^n$ of the form $\pm e_i \pm e_j$ with $i \neq j$. We may take $\Delta = \{e_1 - e_2, e_2 - e_3, \dots, e_{n-1} - e_n, e_{n-1} + e_n \}$. Let $\Phi_0'$ be the standard extremal root subsystem obtained by removing $e_1 - e_2$ from $\Delta$. Then $\Phi_0' = \Phi \cap \{\langle e_1, x \rangle = 0\}$. All other extremal root subsystems are of the form $\Phi_0=w \Phi'_0$; each such one can be expressed as $\Phi\cap \{\langle e_j, x \rangle = 0\}$ for some $j$.

 Suppose $\Phi_0= \Phi\cap \{\langle e_j, x \rangle = 0\}$. If $\Psi$ is obtained by deleting one of the first $n-2$ simple roots, then we proceed as before by looking at the intersection of $\Psi$ with $V_l^n$ and $W_l^n$ (now $\Psi_2$ is of type $D_{n-l}$). If $\Psi$ is obtained by deleting the last simple root, $\Psi$ is just the usual $A_{n-1}$ in $\mathbb{R}^n$ and we are back in the first case. Finally, if $\Psi$ is obtained by removing the second to last root in $\Delta$, then $\Psi$ is also of type $A_{n-1}$ but with basis $e_1-e_2,\ldots, e_{n-2}-e_{n-1}, e_{n-1}+e_n$. But since $\{ \langle e_n, x \rangle = 0 \}= \{ \langle -e_n, x \rangle = 0 \}$, we are again back in the first case.
    
   %\textcolor{gray}{
   % If $j = 1$, then $e_1 - e_2$ is the only simple root not orthogonal to $e_1$. If $2 \leq j \leq n-2$, then there are exactly two simple roots not orthogonal to $e_j$. If $j = n-1$, there are three simple roots not orthogonal to $e_j$, namely $e_{n-2} - e_{n-1}, e_{n-1} - e_n, e_{n-1} + e_n$. If $j = n$, there are two simple roots not orthogonal to $e_j$: $e_{n-1} - e_n, e_{n-1} + e_n$. By the same style of analysis as before, it is clear that if $1 \leq j \leq n-2$, then we are done.} 
    
    %\textcolor{gray}{
    %Now suppose $j = n$. Suppose first that $\Psi$ does not contain the simple root $e_{n-2} - e_{n-1}$. Since $\Psi$ is irreducible and is not contained in $\Phi_0$, it follows that $\Psi$ is the singleton set consisting of either $e_{n-1} - e_n$ or $e_{n-1} + e_n$. This is the degenerate setting referred to earlier, in which \eqref{eqn_root_ineq} reduces to $0 + 1 \geq \frac{1}{2} \cdot 2$. Otherwise, $\Psi$ contains $e_{n-2} - e_{n-1}$, in which case $\Psi$ contains one or both of $e_{n-1} - e_n$ and $e_{n-1} + e_n$. If $\Psi$ contains only one of these extremal roots, then $\Psi$ is of type $A_\ell$, and the same style of analysis as before shows that $\Psi \cap \Phi_0$ is extremal in $\Psi$. Otherwise $\Psi$ contains both $e_{n-1} - e_n$ and $e_{n-1} + e_n$. In this case $\Psi$ is of type $D_\ell$ with corresponding simple roots $e_k - e_{k+1}, \dots, e_{n-2} - e_{n-1}, e_{n-1} - e_n, e_{n-1} + e_n$ (with $\ell = n-k+1$). We can instead choose the basis $e_n - e_k, e_k - e_{k+1}, \dots, e_{n-3} - e_{n-2}, e_{n-2} - e_{n-1}, e_{n-2} + e_{n-1}$. It is then clear that $\Psi\cap \Phi_0$ is extremal in $\Psi$ (i.e., it is of type $D_{\ell-1}$, with the convention that $D_2 = A_1 \times A_1$). }

%\textcolor{gray}{
  %  Finally suppose $j = n-1$. Suppose first that $\Psi$ does not contain the simple root $e_{n-2} - e_{n-1}$. Then, like the $j = n$ case, $\Psi$ is the singleton set consisting of either $e_{n-1} - e_n$ or $e_{n-1} + e_n$, in which case we conclude as before. Otherwise, $\Psi$ contains $e_{n-2} - e_{n-1}$, in which case $\Psi$ either contains exactly one of $e_{n-1} - e_n$ and $e_{n-1} + e_n$ or contains both. In the former case, we are again back in the type $A$ case. Otherwise, $\Psi$ contains both $e_{n-1} - e_n$ and $e_{n-1} + e_n$ and is of type $D_\ell$ with corresponding basis $e_k - e_{k+1}, \dots, e_{n-2} - e_{n-1}, e_{n-1} - e_n, e_{n-1} + e_n$. We instead take the basis $e_{n-1} - e_{k+1}, e_{k+1} - e_{k+2}, \dots, e_{n-2} - e_k, e_k - e_n, e_k + e_n$. It is then clear that $\Psi\cap \Phi_0$ is extremal in $\Psi$ (again with the convention that $D_2 = A_1 \times A_1$). 
    %}
\end{proof}

\end{document}
